\documentclass{hipatia}
\usepackage{lipsum}
%Use \DeclareMathOperator para definir novos
% operadores para o modo matemático
\DeclareMathOperator{\sen}{sen}

%Evite numerar teoremas
%Prefira nomeá-los
%Use os ambientes abaixo
\newtheorem*{theorem*}{Teorema}
\newtheorem*{lemma*}{Lema}

% Evite títulos muito longos
\title{Uma Nova Demonstração do\\
Teorema de Pitágoras}
% Se for necessário diminuir a fonte do título 
% para caber no quadro, use
% \title{ \fontsize{28}{28}\selectfont Uma Nova Demonstração do\\  \fontsize{28}{28}\selectfont Teorema de Pitágoras}

% O Subtítulo é o nome da seção da revista
% Deve ser uma palavra de origem grega
\subtitle{Teorema}
\author{Hipátia de Alexandria e Carl F.  Gauss}
% A data não é necessária
%\date{October 2023}
% Não se preocupe com a numeração
% das páginas ou com o número da edição

\begin{document}

\maketitle

\section{Introdução}

\lipsum[1-3]

\section{Desenvolvimento}
\lipsum[4-5]

\begin{theorem*}
Em um triângulo retângulo, 
a soma dos quadrados dos catetos é igual ao quadrado da hipotenusa.
\end{theorem*}

\begin{figure}[htb!]
\includegraphics[width=8cm]{Pythagorean_proof.png}
\caption{Ideia geométrica da demonstração.}   
\end{figure}

\lipsum[6-7]

\section{Conclusão}

\lipsum[8-9]



\nocite{*}
\vfill

\bibliography{main}


\pagebreak 

% Mini bios 
% Seja informal e divertido
% Prefira fotos com fundo branco
\begin{wrapfigure}{L}{1.7cm}
	\centering
	\includegraphics[width=2cm]{Hypatia_portrait.png}
\end{wrapfigure}\noindent
Hipátia de Alexandria é uma filósofa neoplatônica grega do Egito Romano. É a primeira mulher documentada como tendo sido matemática. Como chefe da escola platônica em Alexandria, também leciona filosofia e astronomia.

\vspace{1cm}
\begin{wrapfigure}{L}{1.7cm}
	\vspace{-10pt}
	\centering
	\includegraphics[width=2cm]{Carl-Friedrich-Gauss-engraving.jpg}
\end{wrapfigure}\noindent
Johann Carl Friedrich Gauss é um matemático, astrônomo e físico alemão que trabalha em diversas áreas da ciência, dentre elas a teoria dos números, estatística, análise matemática, geometria diferencial, geodésia, geofísica, eletroestática, astronomia e óptica.

Alguns se referem a ele como \emph{princeps mathematicorum} (em latim: ``o príncipe da matemática'' ou ``o mais notável dos matemáticos'').

\end{document}
